\documentclass[11pt]{article}

\usepackage{amsmath,amssymb}
\usepackage{xurl}
\usepackage[colorlinks=true,linkcolor=blue,citecolor=blue,urlcolor=blue]{hyperref}

% Shared commands for notes in docs/paper-gaps/.

\DeclareUnicodeCharacter{2080}{\ifmmode{}_{0}\else\textsubscript{0}\fi}
\DeclareUnicodeCharacter{2081}{\ifmmode{}_{1}\else\textsubscript{1}\fi}
\DeclareUnicodeCharacter{2082}{\ifmmode{}_{2}\else\textsubscript{2}\fi}
\DeclareUnicodeCharacter{2099}{\ifmmode{}_{n}\else\textsubscript{n}\fi}

\newcommand{\Lean}{\textsc{Lean}}
\ExplSyntaxOn
\NewDocumentCommand{\leanid}{m}{
  \ifmmode
    \textsf{\detokenize{#1}}
  \else
    \group_begin:
    \urlstyle{sf}
    \tl_set:Nn \l_tmpa_tl {#1}
    \tl_replace_all:Nnn \l_tmpa_tl {\_} {_}
    \exp_args:NV \path \l_tmpa_tl
    \group_end:
  \fi
}
\ExplSyntaxOff
\providecommand{\tr}{\operatorname{tr}}
\newcommand{\tnleanIssueBase}{https://github.com/LionSR/TNLean/issues/}
\newcommand{\tnleanPrBase}{https://github.com/LionSR/TNLean/pull/}
\newcommand{\tnleanDocBase}{https://sirui-lu.com/TNLean/docs/}
\newcommand{\ghissue}[1]{\href{\tnleanIssueBase#1}{GitHub issue \##1}}
\newcommand{\ghpr}[1]{\href{\tnleanPrBase#1}{PR \##1}}
\newcommand{\doclink}[1]{\href{\tnleanDocBase#1.html}{\path{#1}}}

% Dirac notation and the identity operator, as in blueprint/src/macros/common.tex.
% \providecommand keeps note-local definitions authoritative.
\usepackage{dsfont}
\providecommand{\ket}[1]{|#1\rangle}
\providecommand{\bra}[1]{\langle #1|}
\providecommand{\braket}[2]{\langle #1 | #2 \rangle}
\providecommand{\ketbra}[2]{|#1\rangle\!\langle #2|}
\providecommand{\Id}{\mathds{1}}
\providecommand{\one}{\mathds{1}}
\providecommand{\C}{\mathbb{C}}
\providecommand{\MN}[1]{M_{#1}(\C)}
\providecommand{\MD}{M_D(\C)}

% External referencing for paper sources.  The build helper writes sanitized
% auxiliary files under build/paper-aux/ and excludes duplicated source labels.
\usepackage{xr}

% Import a generated paper auxiliary file with a caller-supplied label prefix.
% The candidate paths support builds from docs/paper-gaps/, the repository root,
% and build/paper-aux/ itself.
\newcommand{\importpaperaux}[2]{%
  \IfFileExists{../../build/paper-aux/#2.aux}{%
    \externaldocument[#1]{../../build/paper-aux/#2}%
  }{%
    \IfFileExists{build/paper-aux/#2.aux}{%
      \externaldocument[#1]{build/paper-aux/#2}%
    }{%
      \IfFileExists{#2.aux}{%
        \externaldocument[#1]{#2}%
      }{}%
    }%
  }%
}

% General external reference macros
\newcommand{\paperref}[2]{\ref{#1-#2}}
\newcommand{\papereqref}[2]{(\ref{#1-#2})}

% CPSV16 source (1606.00608 / MPDO-22-12-17-2.tex) external document and shortcuts
\importpaperaux{cpsv16-}{1606.00608/MPDO-22-12-17-2-tnlean}
\newcommand{\cpsvref}[1]{\paperref{cpsv16}{#1}}
\newcommand{\cpsveqref}[1]{\papereqref{cpsv16}{#1}}

% Verdict marker: \gapnote{<kind>}{<status>}. Kind is the policy
% classification (clarification, local-correction, scope-restriction,
% unfaithful, false-source, open-gap); status is open, wip (elimination
% underway), resolved, or historical. Machine-read by the site index;
% prints a verdict line.
\newcommand{\gapnote}[2]{\par\noindent\textbf{Verdict:} #1 (#2).\par}


\title{Star-Closure Correction for the GNVW Support-Algebra Lemma}
\author{}
\date{2026-08-24}

\begin{document}
\maketitle
\gapnote{local-correction}{resolved}

\section{The printed statement and the necessary correction}

The printed statement of arXiv:0910.3675v2, Section~7.1, Lemma~7 says that
$\mathcal A\subseteq\mathcal B_1\otimes\mathcal B_2$ is a subalgebra
\cite{Gross2012IndexTheory}.  It then defines
$\operatorname{Spp}(\mathcal A,\mathcal B_1)$ as the smallest
$C^*$-subalgebra generated by the first-factor coefficients and asserts
\begin{align}
    \operatorname{Spp}(\mathcal A,\mathcal B_1)'
    &=\{b\in\mathcal B_1\mid b\otimes\Id\in\mathcal A'\}.
    \label{eq:gnvw_star_commutant}
\end{align}
The statement and leastness argument occupy lines 1181--1191 of the v2
source, coefficient generation is at lines 1188--1191, and the commutant
argument occupies lines 1192--1208.  The lemma has source label
\path{lem:spp}.  The published DOI
\href{https://doi.org/10.1007/s00220-012-1423-1}
{10.1007/s00220-012-1423-1} and arXiv:0910.3675v2 are the
third-party-verifiable source anchors.

The coefficient-generation and leastness assertions remain valid for an
arbitrary algebra input: coefficient expansion proves containment, and the
dual-basis coefficient maps prove leastness.  The commutant identity
(\ref{eq:gnvw_star_commutant}), however, requires the input to be closed under
the involution.  For a concrete counterexample, take
$\mathcal B_2=\mathbb C$ and
\begin{align}
    \mathcal A&=\operatorname{span}\{\Id,E_{12}\}
      \subset M_2(\mathbb C).
\end{align}
This is a unital algebra but is not closed under adjoints.  Its coefficient
set generates $M_2(\mathbb C)$ as a $C^*$-algebra, since it contains
$E_{12}$ and hence also $E_{21}$.  Consequently the left side of
(\ref{eq:gnvw_star_commutant}) is $\mathbb C\Id$.  The right side is instead
$\mathcal A'=\operatorname{span}\{\Id,E_{12}\}$.  Thus the printed
commutant clause is false for an arbitrary non-star-closed algebra.

To state all clauses as one package, the local correction is to assume that
$\mathcal A$ is a $C^*$-subalgebra, or equivalently in the finite-dimensional
formal setting, a star-subalgebra.  Then the adjoint of every coefficient is
the oppositely indexed coefficient of an adjoint element of $\mathcal A$,
which supplies the involution closure needed in the commutant argument.  This
correction does not assert that coefficient generation or leastness by itself
needs star closure.

\section{Conditional adjoint closure in the applications}

Schumacher and Werner distinguish the ordinary algebraic construction from
the closure property needed in the application to observables
\cite{SchumacherWerner2004QCA}.  They define the support algebra as the
ordinary algebra generated by the coefficient support at lines 1161--1165,
and deduce at lines 1170--1171 that it is closed under adjoints because the
input observable algebra is closed under adjoints.  This is a conditional
adjoint-closure argument; it does not make ``subalgebra'' synonymous with
``$C^*$-subalgebra.''

The GNVW QCA application supplies the corrected hypothesis: at
arXiv:0910.3675v2, lines 1251--1265, the input is the image of a local
observable algebra under a $*$-automorphism.  The MPU application has the
same form.  The map $\omega$ preserves the local star-algebra operations at
lines 2300--2306 of arXiv:1703.09188 \cite{Cirac2017MPU}, and the support
algebra in lines 2313--2334 is taken from
\begin{align}
    \omega(\mathcal A_{2x}\otimes\mathcal A_{2x+1}).
\end{align}
It is therefore the image of a local full matrix star-algebra under a
star-automorphism and is itself star closed.

\section{Formal resolution and remaining scope}

The formal support-algebra API takes its input as a star-subalgebra.  This is
the necessary corrected hypothesis above, and both intended QCA/MPU
applications supply it.  The local correction is therefore resolved: there
is no obligation to formalize the false non-star-closed commutant
characterization.

The distinct restriction to full complex matrix factors remains open and is
recorded in
\path{docs/paper-gaps/gnvw12_support_algebra_full_matrix_scope.tex}.  Removing
that restriction requires a support-algebra construction for arbitrary
finite-dimensional $C^*$-algebra factors; it does not require weakening the
corrected star-closed input.

\bibliographystyle{alpha}
\bibliography{references}

\end{document}
